\chapter{Ingresos energéticos}

\marginal{capítulo con una ausencia declarada}
Este capítulo se escribe sabiendo que le falta una pieza. La trayectoria de la
tasa del derecho por la utilidad compartida, que es el instrumento por el que la
renta petrolera llega al presupuesto, no está en ninguna pieza del paquete y
tampoco en la exposición de motivos del proyecto de presupuesto, que no existe
para este ejercicio. Lo que sigue se sostiene con la Ley de Ingresos y el marco
macroeconómico; la pregunta sobre el instrumento se queda sin responder y se dice
aquí y no en una nota al pie.

\section{Las cuatro vías}

El sector energético entra al artículo primero por cuatro puertas distintas y
conviene no sumarlas sin decirlo, porque no son el mismo objeto: dos son ingreso
propio de empresas, una es renta petrolera transferida y una es un impuesto.

\begin{table}[htbp]
\centering
\caption{El sector energético en el artículo 1o., por vía de entrada (mdp)}
\label{tab:C3_1}
\small
\begin{tabular}{lSSSSS}
\toprule
{Concepto} & {2024} & {2025} & {Dif. nominal} & {Dif. real} & {Var. real \%} \\
\midrule
{Transferencias del Fondo Mexicano del Petroleo (ordinarias)} & \num{277774.3} & \num{279766.8} & \num{1992.5} & \num{-9951.8} & \num{-3.4} \\
{Ingreso propio de Petroleos Mexicanos} & \num{769805.6} & \num{860868.2} & \num{91062.6} & \num{57961.0} & \num{7.2} \\
{Ingreso propio de la Comision Federal de Electricidad} & \num{446951.3} & \num{539145.6} & \num{92194.3} & \num{72975.4} & \num{15.7} \\
{Impuesto por la actividad de exploracion y extraccion de hidrocarburos} & \num{7811.0} & \num{7140.1} & \num{-670.9} & \num{-1006.8} & \num{-12.4} \\
{SUMA de las cuatro vias} & \num{1502342.2} & \num{1686920.7} & \num{184578.5} & \num{119977.8} & \num{7.7} \\
\bottomrule
\end{tabular}
\par\vspace{2pt}\footnotesize Fuente: elaboración propia del ITED con la Ley de Ingresos aprobada de 2024 y 2025. Las diferencias en millones de pesos se reportan dos veces, nominal y real; la real está en pesos de 2025 con el deflactor de 4.3\% que declara el CGPE.
\end{table}


Las transferencias del Fondo Mexicano del Petróleo pasan de 277,774.3 a
\textbf{279,766.8 millones de pesos}, un aumento nominal de 0.7\% que en términos
reales es una caída de 3.4\%. El ingreso propio de Petróleos Mexicanos sube de
769,805.6 a \textbf{860,868.2}, y el de la Comisión Federal de Electricidad de
446,951.3 a \textbf{539,145.6}. El impuesto por la actividad de exploración y
extracción baja de 7,811.0 a 7,140.1.

\marginal{la renta cae, el ingreso propio sube}
El contraste es el hallazgo del capítulo. \textbf{La renta petrolera que llega al
erario por el Fondo cae en términos reales, mientras el ingreso propio de las dos
empresas sube con fuerza.} Con un precio de la mezcla que baja de 70.7 a 57.8
dólares por barril y una plataforma de exportación que baja de 900 a 892 miles de
barriles diarios, la caída de la renta transferida es aritmética. El aumento del
ingreso propio, en cambio, no se explica por el precio, y el paquete no dice de
dónde sale.

\section{Qué se puede afirmar y qué no}

Se puede afirmar que el sector aporta más en total, 1,686,920.7 millones contra
1,502,342.2, y que la composición se desplaza del erario hacia las empresas.

\textbf{No se puede afirmar} cuánto de ese desplazamiento es política deliberada
sobre la tasa del derecho y cuánto es resultado del precio y del volumen. Separar
las dos cosas exige la trayectoria de la tasa, que en ejercicios anteriores esta
casa reconstruyó desde el Diario Oficial y que aquí no se reconstruyó. Es la
ausencia declarada del capítulo.

\clearpage
